
\section{Minimum-Weight Polygon Triangulation}\label{dp: sec: MWT}\doHeader{Min-Weight Triangulation}

\begin{mylist}
\item[\bf input:]
  A sequence $p[1..n]$ of points in the plane,
  forming (in clockwise order) the vertices
  $p[1], p[2], \ldots, p[n]$ of a convex polygon.

\item[\bf output:]
  The minimum weight of any \emph{triangulation} of $p$.
  A triangulation of $p$ is a set of \emph{chords}
  (diagonals---straight-line segments connecting pairs of vertices)
  that subdivides the convex polygon into triangular faces.
  The weight of the triangulation is the sum of the lengths of the chords.
  We include the boundary segments (with endpoints $p_i, p_{i+1}$)
  as chords.
\end{mylist}

\paragraph{subproblems and recurrence relation.}
Fix any input instance $p[1..n]$.
For each pair $(i, j)$ with $1 \le i < j\le n$, and $j-i \ge 1$,
define $T(i, j)$ to be the optimal cost for the subproblem
$p[i..j]$ formed by the $j-i+1 \ge 2$ points $p[i], p[i+1], \ldots, p[j]$.
The final answer is $T(1, n)$.
The recurrence relation is
\[
  T(i, j) = \begin{cases}
    d(p[i], p[j]) & \text{if } j=i+1, \\
    \min\Big\{ T(i, k) + T(k, j) + d(p[i], p[j]) ~\big|~ k\in \{i+1, i+2, \ldots, j-1\} \Big\} & \text{if } j > i+1.
  \end{cases}
\]
Above, $d(p[i], p[j])$ is the Euclidean distance between points $p[i] = (x_i, y_i)$ and $p[j] = (x_j ,y_j)$,
that is, $d(p[i], p[j]) = \sqrt{(x_i-x_j)^2 + (y_i-y_j)^2}$.

\paragraph{correctness.}
The intuition is as follows.
Each triangulation of $p[i..j]$ can be obtained as follows
for some $k\in\{i+1,\ldots,j-1\}$:
take any triangulation of $p[i..k]$,
and any triangulation of $p[k..j]$,
and add the chord with endpoints $p[i]$ and $p[j]$.
(The latter chord must be in the triangulation,
and point $p[k]$ is the third vertex on the triangle
that has that chord as one side.
Draw a picture and consider some examples for more intuition.
Make sure to consider the boundary cases when $k\in\{i+1, j-1\}$.)

\paragraph{time.}
There are $\Theta(n^2)$ subproblems.
For each, the right-hand side of the recurrence can be evaluated in time $O(n)$
(assuming the smaller subproblems are already solved).
So the algorithm (iterative, or recursive and memoized)
takes time $\Theta(n^3)$.

\begin{theorem}
  There is an $O(n^3)$-time algorithm for \prob{Min-Weight Polygon Triangulation}.
\end{theorem}

\paragraph{does not reduce to Shortest Paths!}
The standard trick of creating a graph with one vertex per subproblem $T(i,j)$
doesn't work, because, in the right-hand side of the recurrence for $T(i,j)$,
for each $k$, the term includes \emph{two} subproblems $T(i,k)$ and $T(k,j)$.

\pythonCode{dynamic_programming/min_wt_triangulation.py}

% \paragraph{External sources on \prob{Min-Weight Triangulation}}

% \begin{itemize}

% \item DPV Problem 6.12.
  
% \end{itemize}

%%% Local Variables:
%%% mode: latex
%%% TeX-master: "dynamic_programming"
%%% End:
